\title{Large Language Models Can Automatically Engineer Features for Few-Shot Tabular Learning}

\begin{abstract}
Large Language Models (LLMs) exhibit exceptional capabilities in addressing complex and previously unencountered reasoning tasks, positioning them as highly promising tools for tabular learning, which is essential for a variety of practical applications. In this study, we introduce an innovative in-context learning approach, termed \textsf{FeatLLM}, where LLMs are leveraged as feature designers to construct input datasets specifically tailored for superior tabular prediction performance. The features generated through this process are utilized by a straightforward downstream machine learning algorithm—such as linear regression—to estimate class probabilities, facilitating robust few-shot learning outcomes. Notably, the \textsf{FeatLLM} framework relies solely on the basic predictive model and the extracted features during inference, avoiding the necessity of querying the LLM for every individual data point in real time. Furthermore, this approach only requires API-level access to LLMs and circumvents the constraints posed by prompt length limitations. Across diverse tabular datasets from multiple sectors, empirical results show that \textsf{FeatLLM} produces high-caliber decision rules, yielding substantially better results (with an average improvement of 10\%) when compared to existing methods such as TabLLM and STUNT.
\end{abstract}

\section{Introduction}
Large Language Models (LLMs) have recently demonstrated impressive proficiency in generating contextually appropriate outputs for novel tasks, all without the necessity of task-specific fine-tuning~\cite{creswell2022selection,imani2023mathprompter,taylor2022galactica}. Their extensive training on large-scale textual datasets equips them with deep prior knowledge, effectively reducing reliance on domain experts in numerous fields~\cite{Genomes2021,singhal2023large,wilcox2003role}, thus advancing the automation of various functions ranging from dialogue processing~\cite{finch2023leveraging} to software maintenance~\cite{fan2023automated}. By embedding a task statement and limited examples in prompt formulations, LLMs can grasp contextual nuances, allowing them to prioritize salient features and instances~\cite{brown2020language,zhao2023survey}. This illustrates the potential for LLMs to serve as proficient feature engineers in data analysis.

Recently, the exploitation of LLMs’ background knowledge and reasoning aptitude has expanded into the sphere of tabular learning. For example, domain-specific knowledge—such as the correlation between advanced age and increased disease risk—can inform predictive healthcare applications. Existing approaches typically define tasks and feature metadata in natural language, serialize tabular data, and submit this content to an LLM~\cite{dinh2022lift,wang2023anypredict}. This process is commonly followed by employing either in-context learning paradigms~\cite{nam2023semi} or adopting techniques for efficient parameter fine-tuning~\cite{dinh2022lift,hegselmann2023tabllm}. Empirical findings confirm that the prior knowledge harnessed from LLMs substantially improves performance, especially surpassing conventional tabular learning baselines in settings with limited labeled examples~\cite{hegselmann2023tabllm}.

Existing approaches to tabular learning utilizing LLMs exhibit several notable shortcomings. For direct prediction tasks, these methods necessitate executing at least one LLM inference for every individual sample, resulting in considerable computational overhead. Additionally, achieving strong predictive performance often requires fine-tuning the LLMs, which is impractical for models that lack publicly available training code or tuning APIs. This issue is exacerbated by the prevalence of state-of-the-art LLMs that are only accessible via restricted inference APIs~\cite{anil2023palm,openai2023gpt4}. Furthermore, most current strategies struggle with extended prompt lengths~\cite{liu2023lost}; when the tabular data contains a large number of features, leading to longer textual inputs, the models frequently become impractical or their performance substantially deteriorates. These obstacles create significant barriers for deploying LLM-based tabular learning systems in practical environments.

To address these issues, our method redefines the use of LLMs beyond end-to-end predictions by positioning them as engines for feature engineering, thereby capitalizing more effectively on their extensive prior knowledge. We present a new in-context learning paradigm, \textsf{FeatLLM}, designed to discern and articulate the ``criteria'' that govern LLM decision-making. For example, in medical prediction tasks, LLMs can identify and produce explicit rules indicating which feature combinations correspond to a diagnosis. These extracted rules are employed to generate novel features, supplanting the original ones used in downstream tasks. This transfer of responsibility from prediction to feature creation results in a more streamlined downstream modeling process; once the relevant features are defined, the necessity for sophisticated machine learning inference models diminishes, which leads to improved latency at the inference stage. Harnessing the in-context learning capabilities of LLMs allows feature extraction without traditional model training, simply by augmenting prompt demonstrations. Additionally, ensemble strategies such as feature bagging~\cite{breiman2001random} can be adopted to overcome the difficulties associated with overly long prompts.

\textsf{FeatLLM} organizes the prompt for creating novel features into two distinct subtasks: (1) comprehending the task and deducing the connection between input features and target outcomes; and (2) leveraging this understanding, along with limited example cases, to formulate discriminative rules capable of separating different classes. Through this systematic reasoning framework, large language models are guided to exclude irrelevant features when developing extraction rules. The resulting rules are executed on the dataset (interpreted as executable code), transforming the data into binary indicators that signal whether each sample satisfies individual rules. Subsequently, a simple model is trained to predict class likelihoods based on these newly constructed features. Further enhancement of robustness is accomplished by ensembling, where feature generation and feature bagging are iterated multiple times. By judiciously tuning the number of features and samples for bagging, issues stemming from overly large prompts are addressed. With no dependency on additional model training, \textsf{FeatLLM} maintains low inference demands, enabling LLMs to be applied effectively to diverse practical tasks.

Experiments with \textsf{FeatLLM} span 13 tabular datasets under limited data scenarios, demonstrating high and consistent performance. The results indicate that LLMs are adept at harnessing both pre-existing and empirical knowledge, producing useful features of substantial quality. The proposed method surpasses state-of-the-art few-shot learning approaches under various evaluation protocols. The implementation can be found in an anonymized repository at~\texttt{https://github.com/Sungwon-Han/FeatLLM}.

\section{Related Work}

\subsection{Few-Shot Learning with Tabular Data} Few-shot learning has made it possible to derive transferable and robust representations from datasets, even when labeled instances are scarce~\cite{chen2018closer}. Although this area was originally developed with an emphasis on image datasets, its applicability has broadened to include modalities such as text, audio, and more recently, tabular formats~\cite{majumder2022few,nam2023stunt,schick2021exploiting}. The reliance on limited labeled examples elevates the probability of the model picking up superficial correlations~\cite{bartlett2021deep}. To address this, researchers have explored strategies that supplement the learning process with extra information sources or inject domain-specific priors to guide inductive bias. Approaches include exploiting large collections of unlabeled tabular data paired with augmentation techniques to facilitate semi-supervised learning~\cite{ucar2021subtab,yoon2020vime}, or employing unsupervised pre-training strategies that are task-independent for more effective model initialization~\cite{somepalli2021saint}. Unsupervised objectives are commonly implemented by introducing masked or perturbed data and applying reconstruction-based losses~\cite{arik2021tabnet,majmundar2022met}, or by leveraging data augmentation in support of contrastive learning frameworks~\cite{bahri2022scarf,chen2020simple}. Nonetheless, due to the intrinsic heterogeneity found in tabular datasets, devising universal augmentation strategies remains difficult, and these techniques have typically not shown pronounced success in few-shot contexts~\cite{nam2023stunt}.

A new trend involves training models on an extensive and varied selection of tabular datasets—including those unrelated to the target domain—and subsequently adopting transfer learning techniques for downstream tasks~\cite{zhang2023generative,levin2022transfer,zhu2023xtab}. One example, TransTab~\cite{wang2022transtab}, leverages transformer architectures to learn semantic connections between columns, using column names in addition to the underlying data. Insights drawn from many tables and column structures have enabled models to perform inference on previously unseen tabular tasks in zero-shot and few-shot regimes. Another method, TabPFN~\cite{hollmann2022tabpfn}, combines multiple distribution types to generate artificial data used for pre-training a Bayesian neural network. The model then processes both training and test data for target tasks directly, foregoing further learning stages. This accumulated prior information from diverse datasets has helped these methods outperform traditional tree-based algorithms, especially within the context of few-shot learning.

\subsection{Language-Interfaced Tabular Learning} Integrating large language models (LLMs)~\cite{anil2023palm,openai2023gpt4} provides another avenue for leveraging prior knowledge in tabular learning. Trained on vast and diverse collections of textual data, LLMs have shown considerable adaptability and generalization to novel tasks, demonstrating widespread applicability across several fields~\cite{brown2020language}. Recent research has explored how to capitalize on the broad knowledge present in LLMs to improve tabular data processing. Common approaches involve transforming tabular data into textual form and including these sequences within LLM input prompts~\cite{dinh2022lift,hegselmann2023tabllm,wang2023anypredict}. By embedding tabular data, relevant feature definitions, and detailed task descriptions within prompts, LLMs can extrapolate associations between features and tasks, enabling inference. Advancements in this area have utilized methods that either adapt models via efficient parameter tuning such as LoRA~\cite{hu2021lora} or IA3~\cite{liu2022few}, or rely on in-context learning, which introduces a handful of example tasks in prompts and foregoes additional model training~\cite{dinh2022lift,hegselmann2023tabllm,nam2023semi,slack2023tablet}.

Although previous approaches have shown efficacy in enhancing few-shot tabular learning, their reliance on conducting inference through an LLM for every instance hinders practicality, particularly in industrial settings. This research shifts away from the traditional black-box method that feeds each sample through the LLM in a fully end-to-end manner for prediction. Instead, the focus is on guiding the LLM to deduce class-specific rules or conditions directly. \textsf{FeatLLM} then translates these generated rules into program code, constructing novel features. This allows predictions to be performed without repeatedly invoking the LLM, capitalizing on in-context learning to extract these rules using background knowledge and available few-shot examples.

\section{Methods}
\textsf{FeatLLM} operates by identifying and retrieving "rules"—specific criteria related to each possible answer class—from the limited set of samples it receives as input, as illustrated in Figure~\ref{fig:main_model}. The LLM outputs these rules, which are subsequently transformed into new binary features for each sample, indicating whether a particular rule is fulfilled. The collection of these binary features is then passed through a dot product with a set of non-negative trainable parameters to approximate each class’s likelihood. As the model is trained, it learns which features most influence the prediction for each class (Section~\ref{Sec:training}). This procedure is performed iteratively using bagging and results aggregated through ensembling. The following sections elaborate on the formal problem setup and further details for each stage.

\vspace*{-0.12in}
\paragraph{Problem formulation.} Consider a tabular dataset comprising $N$ labeled entries, denoted by $\mathcal{D} = \{(\mathbf{x}^i, \mathbf{y}^i)\}_{i=1}^{N}$. This dataset $\mathcal{D}$ contains $d$ input attributes, where each $\mathbf{x}^i$ represents a vector with $d$ dimensions. The associated feature names, expressed in natural language and noted as $F = \{f_j\}_{j=1}^{d}$, are supplemented with separate definitions. Within the classification context, the label $\mathbf{y}^i$ corresponds to one of several predetermined classes. In $k$-shot learning scenarios, just $k$ labeled samples ($k<N$) are drawn at random from the dataset for training purposes.

\subsection{Prompt Design for Extracting Rules}
\label{Sec:prompt_design}
To facilitate precise rule extraction by the LLM, we structure the workflow to resemble the analytical strategies employed by expert practitioners when tackling tabular datasets. Our prompt is organized into three core segments as described below (refer to Fig.~\ref{fig:prompt_for_rule} and Appendix~\ref{sec:example_rule}):

\vspace*{-0.12in}
\paragraph{Basic information description.} This section delivers all fundamental elements required to address the task (highlighted in orange in Fig.~\ref{fig:prompt_for_rule}). It covers the specification of the task together with label information, details regarding the features, and illustrative demonstrations. Task descriptions are phrased as queries (for instance, ``Does this patient's myocardial infarction complications data indicate chronic heart failure? Yes or no?"). Additional task description samples can be found in Appendix~\ref{sec:task_desc}. Feature explanations denote both value types and may provide optional definitions. For categorical features, examples of possible values may also be included. Sample demonstrations entail selecting a few representative training instances, formatting them as text alongside their respective ground-truth labels. The serialization adheres to conventions established in prior literature~\cite{dinh2022lift,hegselmann2023tabllm}:

\begin{align}
&\text{Serialize}(\mathbf{x}^i,\mathbf{y}^i, F) = \nonumber \\ 
&``\ f_1\ \text{is}\ \mathbf{x}^i_1.\ ... \ f_d\  \text{is}\  \mathbf{x}^i_d.\ \ \text{Answer:}\  \mathbf{y}^i\ ”
\end{align}

\vspace*{-0.12in}
\paragraph{Reasoning instruction.} To bolster the reasoning capabilities of the LLM, we supply explicit reasoning instructions (highlighted in blue in Fig.~\ref{fig:prompt_for_rule}). The instruction commences with a preliminary sentence, reminiscent of the chain-of-thought prompting method~\cite{wei2022chain}. Subsequently, the inference by the LLM is structured into two stages. In the initial stage, the LLM is prompted to independently infer the causal links or propensities connecting the features to the task description. This reasoning is carried out absent the reference to provided demonstrations, depending instead on the model's general knowledge or common sense. This encourages the LLM to maximize its utilization of prior knowledge relevant to the problem at hand. In the following stage, the LLM leverages both the demonstration examples and its initial analysis to construct classification rules for each class. Employing this dual-phase reasoning discourages the model from relying on superficial correlations, especially those originating from irrelevant columns, and helps concentrate attention on salient feature attributes. To strike a balance between comprehensiveness and prompt brevity, thus avoiding the risks of underfitting or overly lengthy prompts, a predetermined rule count (specifically, 10 rules)\footnote{Refer to Section~\ref{sec:analysis} for an investigation into the effect of rule quantity.} is enforced per class. Furthermore, during rule generation, the LLM is encouraged to consult the feature type information in the basic description section, thereby mitigating errors such as type mismatches in formulated rules.

\vspace*{-0.12in}
\paragraph{Response instruction.} To streamline the extraction and subsequent application of the derived rules, we introduce targeted instructions outlining response formatting requirements (depicted in yellow in Fig.~\ref{fig:prompt_for_rule}). The prompting details specify the expected response structure for each answer category (i.e., response format) as well as the template each rule should adhere to (i.e., rule format). These directives empower the LLM to adjust its formatting according to the specific feature types encountered. Even absent further instruction, the model is able to interlink rules using logical connectors such as AND or OR. For an example of the resultant output, consult Appendix~\ref{sec:outcome-rule}.

\subsection{Inferring Class Likelihood via Rules}
\label{Sec:training}

\paragraph{Parsing rules for feature generation.} The rules produced in the prior phase serve as the foundation for constructing novel binary features. These features are class-specific and denote whether a particular instance adheres to the corresponding rules (represented as ‘0’ or ‘1’). Such features can subsequently inform predictions regarding a sample’s class membership. Despite enforcing structured formatting during rule creation to facilitate parsing, the LLM-generated rules—articulated in natural language—often introduce syntactic irregularities~\cite{kovavc2023large}, complicating automatic feature extraction. For instance, the model may replace symbols like ``$>$” and ``$<$” with their textual equivalents, such as ``greater than” or ``less than”, leading to unpredictable variations.

To manage the difficulties arising from parsing inconsistent or noisy textual rules, we opt to utilize the LLM itself instead of devising intricate parser code. The LLM excels at comprehending semantic content amid syntactic fluctuation, and it is capable of synthesizing straightforward code given proper guidance, thanks to its extensive exposure to programming corpora. To capitalize on this, our prompt incorporates the function name, specifications of the inputs and outputs, and the derived rules, which are fed directly to the LLM. Illustrative prompts and their respective parsing outputs are provided in Figures~\ref{fig:prompt_for_parsing} and ~\ref{sec:example_parsing}-\ref{sec:example_parse_outcome} in the Appendix. The resultant code is executed with Python’s \texttt{exec()} to convert the rules into actionable features.

\vspace*{-0.12in}
\paragraph{Inferring class likelihood.} Once the binary features from class-specific rules have been established, a straightforward approach to estimate a sample’s class likelihood is to tally the number of satisfied rules per class (i.e., aggregate the corresponding binary features). However, since individual rules and features may hold different predictive value, it is necessary to assess their relative importance, drawing on data from the training set. We propose determining these weights using standard machine learning models, applied individually to each set of binary features per class. For example, a linear model without intercept can be employed. Suppose $\mathbf{z}_k^i$ represents the binary feature vector generated from sample $\mathbf{x}^i$ with respect to class $k$ such that $\mathbf{z}_k^i \in \{0, 1\}^R$, where $R$ denotes the number of rules per class and $c$ refers to the number of classes. The class probability vector $\mathbf{p}^i$ is then calculated as follows:

\begin{align}
\text{logit}_k^i &= \max(\mathbf{w}_k, 0) \cdot \mathbf{z}_k^i, \nonumber \\
\mathbf{p}^i &= \text{Softmax}({[\text{logit}_{1}^i, ..., \text{logit}_{c}^i]}). \label{eq:infer_prob}
\end{align}

In this context, $\mathbf{w}_k$ denotes the adjustable parameters in the linear model and serves as an indicator for each feature's contribution. By constraining these weights to lie within the positive region, we guarantee that features produced by the LLM cannot negatively impact the predicted probability of a class during training. After obtaining the logit vector, we apply the softmax function to transform it into probabilities associated with each class. Training proceeds by minimizing the cross-entropy loss, where $\mathbf{w}_k$ is estimated using a limited set of labeled samples.

\vspace*{-0.12in}
\paragraph{Ensembling with bagging.} The entire procedure is performed multiple times to generate several inference models, which are subsequently aggregated to form the final prediction using an ensemble. For $T$ repetitions, the corresponding trained weights $\{\mathbf{w}[t]\}_{t=1}^T$ are used to calculate the probability vector for each class, $\mathbf{p}^i$, with the average over all $T$ trials serving as the output, as expressed in Eq.~\ref{eq:infer_prob}.

To inject diversity into the rules produced by the LLM for ensembling, we set the temperature parameter sufficiently high during inference. Furthermore, we shuffle the order of few-shot exemplars, since prior findings indicate that this sequencing can alter the LLM's responses~\cite{lu2022fantastically}\footnote{Further discussion on rule diversity is provided in Appendix~\ref{sec:diversity}}. Bagging is utilized to harness this variation by randomly selecting subsets of features or training instances for each trial, a technique particularly valuable as sample or feature counts grow. This ensemble scheme provides two primary benefits. Firstly, erroneous rules from some LLM generations may be offset by correct rules from others, boosting the system's robustness. This phenomenon echoes the notion of LLM self-consistency~\cite{wang2022self}, whereby rules that recur across multiple ensemble members tend to be trustworthy. Secondly, the ensemble mechanism mitigates issues arising from the limited input prompt size permissible for the LLM. As datasets become large and exceed these input limits, bagging ensures that performance remains stable by breaking data into manageable pieces. Although any single set of rules may not incorporate all features, pooling weak learners built from various trials enables the system to minimize the effects of incomplete feature and instance representation.

\section{Experiments}
We assess the effectiveness of \textsf{FeatLLM} across several tabular datasets and carry out a range of investigative analyses to elucidate the workings of our framework. Owing to page limitations, supplementary experimental details, including tests on different LLM architectures and qualitative evaluations, are provided in the Appendix.

\subsection{Performance Evaluation}
\paragraph{Datasets.}
Our empirical study draws upon 13 datasets targeting binary or multi-class classification: (1) Adult~\cite{asuncion2007uci}, which determines whether a person’s annual income surpasses \$50,000; (2) Bank~\cite{moro2014data}, which assesses the likelihood of a client subscribing to a term deposit; (3) Blood~\cite{yeh2009knowledge}, which predicts donor return for future blood donations; (4) Car~\cite{kadra2021well}, which evaluates car quality; (5) Communities~\cite{misc_communities_and_crime_183}, which forecasts regional crime rates; (6) Credit-g~\cite{kadra2021well}, which estimates individual creditworthiness; (7) Diabetes\footnote{\texttt{kaggle.com/datasets/uciml/pima-indians-diabetes-database}}, which predicts diabetes diagnosis; (8) Heart\footnote{\texttt{kaggle.com/datasets/fedesoriano/heart-failure-prediction}}, which identifies coronary artery disease risk; and (9) Myocardial~\cite{misc_myocardial_infarction_complications_579}, which predicts the presence of chronic heart failure.

In addition, our collection incorporates newer datasets and synthetically generated ones that are seldom discussed and omitted from the LLMs' training data: (10) Cultivars, which predicts soybean cultivar grain yield\footnote{\texttt{archive.ics.uci.edu/dataset/913}}; (11) NHANES, which classifies a person’s age bracket\footnote{\texttt{archive.ics.uci.edu/dataset/887}}; (12) Sequence-type, which categorizes a sequence as arithmetic, geometric, Fibonacci, or Collatz; (13) Solution-mix, which determines if the concentration level of a mixture from four solutions surpasses 0.5. Cultivars and NHANES were contributed to the UCI Machine Learning Repository post-September 2023, guaranteeing their exclusion from the pre-training of the LLM API (gpt-3.5-turbo-0613) employed in our framework. Sequence-type and Solution-mix are synthetic datasets that further prevent memorization by the LLM, thus facilitating unbiased evaluation.

There is substantial variation in both dataset size and complexity, which is further outlined in Table~\ref{Tab:basic-desc}.
Each dataset delineates explicit names and explanations for individual features. ‘Communities’ and ‘Myocardial’, notably, contain over 100 attributes, presenting additional difficulties due to limited LLM context window capabilities.

\vspace*{-0.12in}
\paragraph{Baselines.}
We benchmark our framework against nine distinct baselines. The initial three reference traditional tabular modeling: (1) Logistic regression (LogReg), (2) XGBoost~\cite{chen2016xgboost}, and (3) RandomForest~\cite{ho1995random}. Two further baselines depend on availability of unlabeled data: (4) SCARF~\cite{bahri2022scarf}, and (5) STUNT~\cite{nam2023stunt}. Note that in many practical scenarios, amassing large volumes of unlabeled data proves difficult. Another contemporary method is (6) TabPFN~\cite{hollmann2022tabpfn}, which utilizes synthetic, varied-distribution datasets for pre-training. Finally, three baselines employ LLM methodologies: (7) In-context learning (In-context)~\cite{wei2022emergent}, (8) TABLET~\cite{slack2023tablet}, and (9) TabLLM~\cite{hegselmann2023tabllm}. In-context learning involves embedding few-shot examples directly within the prompt, with no parameter adaptation. TABLET augments inference quality by introducing supplementary information via rule sets and prototypes obtained from an external classifier. TabLLM applies parameter-efficient tuning (such as IA3~\cite{liu2022few}), enabling LLM fine-tuning with limited examples.

\vspace*{-0.12in}
\paragraph{Implementation details.}
\textsf{FeatLLM} utilizes GPT-3.5 as its primary LLM, although it is constructed to function independently of any particular LLM architecture (see Appendix~\ref{sec:backbones} for experimental results across various LLMs). The inference temperature for the LLM is configured at 0.5, with the top-p parameter maintained at the default value of 1 in the API settings. For ensemble generation and rule extraction, the respective numbers are set to 20 and 10. Analysis of hyper-parameter effects can be found in Figure~\ref{fig:hyperparameter2} and Appendix~\ref{sec:hyper-parameter}. 

Adam optimizer is employed for the linear model, with a learning rate fixed at 0.01 and trained over 200 epochs. Epoch selection is performed through $k$-fold cross-validation. For baseline replication, GPT-3.5 is applied for in-context learning and T0~\cite{sanh2022multitask} is used for TabLLM in accordance with the original protocol. TabLLM limits model choice to LLMs that provide publicly available checkpoints, thus excluding GPT-3.5 from usage in this baseline. Parameter selection for classic machine learning algorithms (LogReg, XGBoost, RandomForest) is accomplished through grid-search combined with $k$-fold cross-validation, with $k$ set to either 2 or 4 to ensure coverage of every class within the training dataset. Additional implementation specifics can be found in Appendix~\ref{sec:baselines}.

\vspace*{-0.12in}
\paragraph{Results.}
Tables~\ref{tab:main1} and \ref{tab:main2} provide comparative results on multiple datasets, where each experimental setting is conducted three times with varied random splits for the training and test data; we report both the average and standard deviation of the obtained AUC scores. \textsf{FeatLLM} reliably emerges as either the leading method or the runner-up in nearly all cases relative to alternative baseline methods. Remarkably, our method demonstrates strong sample efficiency: leveraging only 4 shots, it attains levels of performance on par with standard tabular baselines that utilize 32 shots, as illustrated in Fig.~\ref{fig:compares}. Although STUNT and TabPFN offer notable improvements for particular datasets, their gains over classical machine learning strategies are generally modest. Further, models leveraging LLMs (such as in-context learning, TABLET, and TabLLM) are unable to process tabular datasets characterized by a large feature space. In contrast, our approach, which incorporates both bagging and ensemble strategies, maintains robust effectiveness even when the datasets involve many features. This performance is consistently satisfactory. It should be emphasized that integrating our bagging and ensemble methodology with existing LLM-based frameworks is theoretically possible; however, this adaptation would necessitate doubling the number of inferences per instance, drastically escalating computational demands to impractical levels.

\subsection{Analysis \& Discussion}
\label{sec:analysis}
\paragraph{Ablation study.} We investigate the influence of key framework components on overall performance via ablation experiments, specifically examining: (1) \textsf{-Tuning}: removing the process of weight optimization within the linear model and instead aggregating the binary feature directly for logit computation; (2) \textsf{-Ensemble}: removing the ensemble component; (3) \textsf{-Description}: excluding feature descriptions; and (4) \textsf{-Reasoning}: omitting the initial reasoning step (Step-1) in the instruction used for rule generation.

Table~\ref{tab:ablation} reports the averaged AUC change attributed to each component’s removal, aggregated across all datasets. Across the board, eliminating any of these mechanisms leads to diminished performance. Notably, the utility of weight tuning grows with increasing shot counts, revealing that with ample data, more precise estimation of rule significance becomes possible, thereby amplifying the impact of tuning. By contrast, when only limited shots are available, feature description and reasoning instructions exert greater influence on outcomes; this highlights the importance of effectively leveraging LLM prior knowledge in data-scarce settings. Finally, our ensemble strategy persistently yields performance boosts regardless of the scenario.

\vspace*{-0.12in}
\paragraph{Training \& inference time comparison.} We assess the efficiency of different approaches in terms of both training and inference duration. The Adult dataset serves as the basis for experimentation, with a training set consisting of 16 shots. All models, with the exception of TabLLM, are executed on a single A100 GPU, whereas TabLLM employs four A100 GPUs to facilitate model parallelism. Table~\ref{tab:cost} presents the comparison of runtime across methods. Importantly, \textsf{FeatLLM} achieves inference speeds on par with traditional machine learning models, reflecting its efficiency. Conversely, LLM-oriented techniques—including In-context learning, TABLET, and TabLLM—require substantially more time for inference. Regarding training, \textsf{FeatLLM} demonstrates training times similar to other LLM-based systems, necessitating merely 30 API calls. This modest requirement is both cost-effective and can be handled concurrently, minimizing resource overhead.

\vspace*{-0.12in}
\paragraph{Quality of features generated by \textsf{FeatLLM}.}
To evaluate the effectiveness of rule-based features generated by \textsf{FeatLLM} for real-world tasks, we implement a straightforward analysis. This involves computing the average AUROC between each newly constructed feature and the corresponding target labels, then determining the proportion of features with AUROC values exceeding 0.5 within every dataset. Table~\ref{tab:quality} compiles these findings, demonstrating that most of the produced rules maintain strong relevance to the target variable and contribute meaningful information toward task resolution (see ``Before tuning"). Furthermore, as the linear model undergoes tuning on the dataset, irrelevant or weakly predictive rules (characterized by learned weights below the specified threshold) are systematically eliminated (see ``After tuning"). The threshold is set to 0.1, guided by the distribution of weight values within the overall histogram.

\vspace*{-0.12in}
\paragraph{Impact of introducing spurious correlations.} To evaluate how well different approaches manage to eliminate irrelevant spurious correlations under low-shot conditions, we performed an experimental assessment. Using the Heart dataset as the main task, we artificially appended columns from the Adult dataset—which have no meaningful association with the Heart prediction objective. We assessed model performance as the number of non-informative columns was incrementally increased. For the sake of consistency, each method was provided with exactly 8 labeled samples, disregarding unlabeled data, meaning methods such as SCARF or STUNT were omitted from this analysis.

In Figure~\ref{fig:spurious}, we present the effect of added spurious correlations on model performance. Among all tested methods, \textsf{FeatLLM} demonstrated the smallest reduction in accuracy attributable to spurious features. This outcome is attributable to our framework’s enforcement of task-feature linkage—drawn from prior knowledge—during rule extraction, thereby ensuring rules aren’t created from unrelated columns and enhancing model resilience. Quantitatively, rules referencing noisy columns contributed only 1–2\% to the overall rule set. Conversely, we note that \textsf{FeatLLM} is not immune to learning misleading correlations, especially when columns from conceptually similar datasets are randomly appended; this can lead to faulty reasoning about causal links between task and features (see Appendix~\ref{sec:spurious-appendix}).

\vspace*{-0.12in}
\paragraph{Hyper-parameter analysis}
\textsf{FeatLLM} relies on two primary hyper-parameters: the quantity of ensembles and the number of rules extracted per class in each LLM query. We systematically evaluated how varying these hyper-parameters influences overall performance. Experimental findings indicate that increasing the ensemble count enhances performance, yet improvements plateau once the ensemble number approaches 20 (refer to Fig.~\ref{fig:appendix-hyperparameter1} in Appendix). For the rule count per class, statistical analysis revealed no significant differences across choices. Balancing prompt size and accuracy, 10 rules per class emerged as optimal (see Fig.~\ref{fig:hyperparameter2}). This observation suggests that adding more rules expands the input space and information, but also amplifies the risk of overfitting in settings with limited data.

\section{Conclusion}
We introduce \textsf{FeatLLM}, which substantially advances LLM-based few-shot tabular learning. Distinct from methods that merely leverage LLMs for pointwise sample predictions, \textsf{FeatLLM} utilizes LLMs for generating predictive criteria, acting as a feature engineering agent. Through its rule extraction mechanism coupled with ensemble bagging, \textsf{FeatLLM} achieves efficient inference, requires only API-level access without additional training, and circumvents restrictions imposed by feature dimensionality. Notably, it demonstrates excellent prediction accuracy and stands as a robust, data-efficient solution for LLM application to practical tabular datasets.

\textsf{FeatLLM} is specialized for few-shot learning situations. Moving forward, we aim to broaden its scope to construct new features in settings with larger sample sizes, explore feature representations that extend beyond rule-based approaches, and enhance interpretability, thereby making it suitable for a wider variety of use cases.

\section{Broader Impact}
The \textsf{FeatLLM} framework is designed to harness LLMs’ prior knowledge for tabular tasks, surpassing standard supervised learning methods especially under data-scarce conditions. By leveraging such prior information, our approach addresses overfitting risks inherent to limited sample training and boosts data efficiency. \textsf{FeatLLM} is especially advantageous for practitioners in fields like finance and healthcare, where data annotation is burdensome and relevant domain knowledge within LLMs may be leveraged—since these models are often pre-trained on pertinent textual resources. As a pivotal direction for future research and broader deployment, we recognize the necessity to rigorously assess the societal bias and misinformation that may be encoded within LLMs’ prior knowledge, given their potential to shape or even outweigh predictions relative to empirical data.

\end{document}